\section{Solution Design}
\label{sec:design}

\subsection{Interface}
\label{sec:interface}

The code generation is one Acceleo module, \texttt{render.mtl}, in the
\texttt{render} bundle, and one Java class, \texttt{Rendering}, that runs it.
The module's input is the target metamodel: its main template takes one
\texttt{ResolvedDeployment}. Its output is a tree of files under a root
directory, each at the path the resolved model assigns it. The module reads
nothing else: no file, no environment, no clock.

\texttt{Rendering.render} takes a list of resolved models and a root
directory. It copies the models into one resource, so the caller's models are
left as they were, and so a query can see every model rendered together
(Section~\ref{sec:sharedindex}). It then loads the module through Acceleo's
standalone API and generates. Any diagnostic Acceleo reports, whether a parse
error in the module or an evaluation error in a template, fails the run with
Acceleo's own message, so a half-written tree is never mistaken for a whole
one. Listing~\ref{lst:runner} shows the core of the runner.

\begin{lstlisting}[language=javarunner, style=report, caption={The Runner, Abridged}, label={lst:runner}]
ResourceSet resources = new ResourceSetImpl();
Resource model = new XMIResourceImpl(URI.createURI(MODEL));
resources.getResources().add(model);
model.getContents().addAll(EcoreUtil.copyAll(deployments));
ClassLoaderQualifiedNameResolver resolver = new ClassLoaderQualifiedNameResolver(
        Rendering.class.getClassLoader(), resources.getPackageRegistry(), "::");
IQualifiedNameQueryEnvironment environment =
        AcceleoUtil.newAcceleoQueryEnvironment(Map.of(), resolver, resources, false);
// A template's parameter type resolves only against the packages its environment knows.
environment.registerEPackage(ResolvedDeploymentPackage.eINSTANCE);
AcceleoEvaluator evaluator = new AcceleoEvaluator(environment.getLookupEngine(), "\n");
resolver.addLoader(new ModuleLoader(new AcceleoParser(), evaluator));
AcceleoUtil.generate(evaluator, environment, (Module) resolver.resolve("render"), model,
        new DefaultGenerationStrategy(resources.getURIConverter(), new DefaultWriterFactory()),
        URI.createFileURI(root.toAbsolutePath() + "/"), null, new BasicMonitor());
if (evaluator.getGenerationResult().getDiagnostic().getSeverity() != Diagnostic.OK)
    throw new IllegalStateException("the rendering did not complete: " + ...);
\end{lstlisting}

The pipeline's entry point, \texttt{Pipeline.render}, puts the three tasks in
one call: it parses and checks every document, lowers and resolves each project
named, and hands the resolved models to \texttt{Rendering.render}. The build
calls it for the union of the \texttt{minimal} and \texttt{data} projects and
leaves the tree under \path{bundles/cli/target/parity/rendered/}, where the
parity suite compares it with the oracles (Section~\ref{sec:strategy}).

\subsection{The Shape of the Module}
\label{sec:shape}

The module has 35 templates and 16 queries. Figure~\ref{fig:calls} shows how
they call each other.

\begin{itemize}
  \item \textbf{The main template} walks the resolved model's generated-file
    entries, the \emph{deliverables}, and opens one file per entry at its path
    (Listing~\ref{lst:main}).
  \item \textbf{\texttt{document}} writes the one header line every YAML file
    carries, or writes nothing before a JSON document, which cannot carry a
    comment. \textbf{\texttt{yaml}} then dispatches on the kind of file.
  \item \textbf{One \texttt{objects} template per kind of file} spells that
    file's format: sixteen kinds, from a namespace to a backup job.
    Appendix~\ref{app:templates} lists them, with the objects each writes.
  \item \textbf{Fragments} are what several objects share: the metadata block,
    the fixed label set, a network peer, a probe target, a container, an
    environment list. Each is written once, unindented, and every caller
    places it at its own depth (Section~\ref{sec:margins}).
  \item \textbf{Queries} hold every rule about spelling: durations as seconds,
    precedences as priorities, quoting, the name of a volume, the shared index
    (Section~\ref{sec:spellingrules}).
\end{itemize}

\begin{lstlisting}[language=acceleo, style=report, caption={The Main Template and the Header}, label={lst:main}]
[comment @main /]
[template public render(deployment : resolveddeployment::ResolvedDeployment)]
  [for (d : resolveddeployment::Deliverable | deployment.deliverables)]
    [file (d.path, overwrite, 'UTF-8')]
      [d.document()/]
    [/file]
  [/for]
[/template]

[template private document(d : resolveddeployment::Deliverable)]
  [if (d.oclIsKindOf(resolveddeployment::VaultPolicyFile))]
    [d.oclAsType(resolveddeployment::VaultPolicyFile).json()/]
  [elseif (d.oclIsKindOf(resolveddeployment::VaultRoleFile))]
    [d.oclAsType(resolveddeployment::VaultRoleFile).json()/]
  [else]
    # GENERATED. Never hand-edit.
    [d.yaml()/]
  [/if]
[/template]
\end{lstlisting}

\begin{figure}[!htbp]
\centering
\begin{tikzpicture}[
  node distance=0.35cm and 0.45cm,
  t/.style={draw, rounded corners=2pt, align=center, font=\scriptsize,
    minimum height=0.6cm, text width=1.9cm, fill=green!8},
  f/.style={t, fill=blue!6},
  q/.style={t, fill=orange!10},
  arr/.style={-{Stealth[length=1.6mm]}}]
\node[t] (main) {render (main)};
\node[t, right=of main] (doc) {document};
\node[t, right=of doc] (yaml) {yaml};
\node[t, above right=0.2cm and 0.45cm of yaml] (json) {json (2)};
\node[t, right=of yaml] (obj) {objects (14)};
\node[f, below=0.6cm of obj] (frag) {metadata, labelSet, peer, container, env, probeTarget, \dots};
\node[q, left=of frag] (query) {seconds, scalar, listed, middlewareName, \dots};
\draw[arr] (main) -- (doc);
\draw[arr] (doc) -- (yaml);
\draw[arr] (doc) |- (json);
\draw[arr] (yaml) -- (obj);
\draw[arr] (obj) -- (frag);
\draw[arr] (obj) -- (query);
\draw[arr] (frag) -- (query);
\end{tikzpicture}
\caption{How the Templates and Queries Call Each Other}
\label{fig:calls}
\end{figure}

\subsection{Whitespace: the Margin Rule}
\label{sec:margins}

YAML encodes structure in indentation, so the whitespace a template writes is
part of the result. Acceleo~4 has two rules about it, and the module's
structure follows from them.

\begin{enumerate}
  \item \textbf{Every line of a block's body carries two spaces of margin,
    which are not written.} The body of a template, a \texttt{for} or an
    \texttt{if} is indented two spaces deeper than its header, and those two
    spaces are removed from the output. A line with less margin is a parse
    error.
  \item \textbf{A call written at a column repeats that column's text before
    every line the call writes.} If a template is called after six spaces, the
    six spaces are written before every line of its output.
\end{enumerate}

The second rule is what makes fragments possible. The label set is one
template that writes five lines starting at column zero
(Listing~\ref{lst:labels}). The metadata block calls it one level under its
own \texttt{labels:} key, and the workload's pod template calls it several
levels deeper. The same five lines therefore appear at every depth YAML needs
without being written twice.

\begin{lstlisting}[language=acceleo, style=report, caption={A Fragment, Placed at Its Caller's Depth}, label={lst:labels}]
[template private metadata(p : resolveddeployment::ResolvedProcess)]
  name: [p.name/]
  namespace: [p.application().namespace/]
  labels:
    [p.labels()/]
[/template]

[template private labelSet(name : String, component : String, application : String)]
  app.kubernetes.io/name: [name/]
  app.kubernetes.io/instance: [name/]
  app.kubernetes.io/part-of: [application/]
  app.kubernetes.io/managed-by: deploy-kit
  app.kubernetes.io/component: [component/]
[/template]
\end{lstlisting}

The same rule has one trap. The repeated text is copied literally, so a call
written after \texttt{- } would start \emph{every} line with a dash, and
turn each key of a container into a list item of its own. A fragment that is a
list item therefore opens with its own \texttt{- } and indents the rest of its
lines two spaces, and its caller places it with spaces only. The container
fragment of Listing~\ref{lst:workload} is written that way. Conditional lines
follow the first rule: an \texttt{if} adds a level of margin, so a line inside
it is indented two spaces more in the module than in the output, as the
\texttt{fsGroup} line shows.

\begin{lstlisting}[language=acceleo, style=report, caption={Part of the Workload Template, and the Container Fragment}, label={lst:workload}]
        securityContext:
          runAsNonRoot: true
          runAsUser: [p.uid/]
          runAsGroup: [p.gid/]
  [comment A volume is written as the image's group, and only a volume is. /]
  [if (p.volumes->notEmpty())]
            fsGroup: [p.gid/]
  [/if]
          seccompProfile:
            type: RuntimeDefault
        containers:
          [p.container()/]

[template private container(p : resolveddeployment::ResolvedProcess)]
  - name: [p.name/]
    image: [p.image/]
  [if (p.surfaces->notEmpty())]
      ports:
    [for (s : resolveddeployment::ResolvedSurface | p.surfaces)]
          - name: [s.name/]
            containerPort: [s.port/]
    [/for]
  [/if]
\end{lstlisting}

\subsection{Spelling Rules}
\label{sec:spellingrules}

Every rule about how a value is written is one query, so a template never
computes anything itself.

\begin{itemize}
  \item \textbf{Durations.} The model records a duration as text, such as
    \texttt{60s}; Kubernetes fields take whole seconds. \texttt{seconds}
    converts seconds, minutes and hours.
  \item \textbf{Quoting.} A string is written plain unless the YAML~1.2 core
    schema would read it as null, a boolean, an integer or a float, in which
    case it is double-quoted, as the production implementation's serializer
    does (Listing~\ref{lst:scalar}). The environment value \texttt{50} is
    therefore written \texttt{"50"}, and \texttt{production} is written plain.
  \item \textbf{Block scalars.} A configuration file's text is written as a
    literal block (\texttt{|}), each line indented under its key, an empty line
    left empty, and its last line break kept the way the serializer keeps it.
  \item \textbf{Names.} A volume for a writable path is named after the path,
    and a file mounted from a configuration file is named after the last
    segment of its source.
  \item \textbf{Priorities and middleware.} A route's precedence becomes a
    proxy priority, \texttt{1000} minus the precedence, and a middleware step
    becomes the name of the middleware object the edge project renders.
\end{itemize}

\begin{lstlisting}[language=acceleo, style=report, caption={Two Spelling Queries}, label={lst:scalar}]
[query private seconds(duration : String) : Integer =
  let amount : Integer = duration.substring(1, duration.size() - 1).toInteger() in
  if duration.endsWith('h') then amount * 3600
  else if duration.endsWith('m') then amount * 60 else amount endif endif /]

[query private scalar(value : String) : String =
  if value.matches('|~|null|Null|NULL|true|True|TRUE|false|False|FALSE|[-+]?[0-9]+|0o[0-7]+|0x[0-9a-fA-F]+|[-+]?([.][0-9]+|[0-9]+([.][0-9]*)?)([eE][-+]?[0-9]+)?|[-+]?[.](inf|Inf|INF)|[.](nan|NaN|NAN)')
  then '"' + value + '"' else value endif /]
\end{lstlisting}

\subsection{The Shared Index}
\label{sec:sharedindex}

The index of an entry point's directory lists the routes of every project the
entry point serves (Section~\ref{sec:shared}). Its template does not read its
own entry's list but calls \texttt{listed}
(Listing~\ref{lst:listed}), which collects the entries of every index file at
the same path in every resolved model rendered together, without duplicates,
in name order. Each model then writes the same file, so writing it twice is
harmless and the order of the models does not matter. A first version read the
other models through the resource that holds them; AQL cannot navigate a
resource, because a resource is not a model element, so the query uses
\texttt{allInstances} instead, which the query environment answers over the
resource set it was given.

\begin{lstlisting}[language=acceleo, style=report, caption={The Index Template and Its Query}, label={lst:listed}]
[template private objects(d : resolveddeployment::IndexFile)]
  ---
  apiVersion: kustomize.config.k8s.io/v1beta1
  kind: Kustomization
  resources:
  [for (r : String | d.listed())]
      - [r/]
  [/for]
[/template]

[query private listed(d : resolveddeployment::IndexFile) : Sequence(String) =
  resolveddeployment::IndexFile.allInstances()->select(i | i.path = d.path).resources
    ->asOrderedSet()->sortedBy(r | r) /]
\end{lstlisting}

\subsection{What the Templates Leave to the Transformation}
\label{sec:boundary}

Because the templates decide nothing, several decisions that look like
formatting are made by the transformation and only read here.

\begin{itemize}
  \item \textbf{Which files exist.} The transformation plans a file only when
    it would hold at least one object, and assigns its path. A template never
    skips a file or invents a path.
  \item \textbf{What each index lists.} The transformation lists, for every
    directory, the files a cluster applies from it, leaving out the network
    policies and the secret store's documents, which nothing applies yet.
  \item \textbf{Which monitor a process gets.} A blue-green process is
    scraped through its pods and any other through its service; the
    transformation records the kind, and the template writes
    \texttt{PodMonitor} or \texttt{ServiceMonitor} from it.
  \item \textbf{Who holds a secret.} The transformation lists, per
    application, every identity that synchronises a secret: each process, and
    the backup identity that holds the backup's credential.
\end{itemize}

\subsection{Running the Code Generation}
\label{sec:running}

The build runs the templates headless: \texttt{./mvnw clean verify} in
\texttt{emf/} runs the render bundle's own tests, then the pipeline, which
writes the rendered tree, then the parity suite, which compares it. In Eclipse
Modeling Tools with the Acceleo~4 SDK, \texttt{render.launch} in the render
bundle runs the same module on the resolved model of \texttt{minimal} that the
build leaves in the parity output, and writes its files under the bundle's
\path{target/render/}.
